\section{构网型电池储能模型}
\label{sec:battery}

电池模型同时描述小时级能量状态和构网运行所需的功率、能量备用。图~\ref{fig:battery-model}给出了主要状态变量和约束关系。
\begin{figure}[H]
\centering
\resizebox{0.88\textwidth}{!}{\begin{tikzpicture}[
  font=\small,>={Stealth[length=2mm]},
  bus/.style={draw=blue!65!black,fill=blue!8,rounded corners=2pt,
    minimum width=30mm,minimum height=13mm,align=center,font=\bfseries},
  box/.style={draw=orange!80!black,fill=orange!10,rounded corners=2pt,
    minimum width=33mm,minimum height=12mm,align=center},
  state/.style={draw=violet!65!black,fill=violet!8,rounded corners=2pt,
    minimum width=34mm,minimum height=12mm,align=center},
  flow/.style={->,thick,draw=gray!80},
  biflow/.style={<->,very thick,draw=gray!80},
  info/.style={->,thick,dashed,draw=blue!65}
]
\node[bus] (bus) at (0,0.7) {公共交流母线};
\node[box] (pcs) at (6,0.7) {双向PCS\\充放电互斥};
\node[box] (cell) at (10,0.7) {电池能量体\\效率/自放电};
\node[state] (soc) at (10,-1.25) {时段末$E_t^B$/SOC\\状态递推与期末规则};
\node[state] (gfm) at (6,2.45) {构网责任$g_t$\\功率/能量备用代理};
\node[state] (audit) at (10,2.45) {REGFM--EMT--HIL\\动态阶段门};
\draw[biflow] (bus) -- node[midway,above]{母线侧$P_t^{ch},P_t^{dis}$} (pcs);
\draw[biflow] (pcs) -- (cell);
\draw[flow] (cell) -- (soc);
\draw[info] (soc) -| (pcs);
\draw[info] (gfm) -- (pcs);
\draw[info] (gfm) -- (audit);
\end{tikzpicture}
}
\caption{构网型电池储能子模型结构}
\label{fig:battery-model}
\sourcetype{本研究绘制。}
\end{figure}

\subsection{能量状态与功率边界}

电池储能承担秒级--分钟级构网支撑与小时级能量搬移。小时模型以母线侧充放电功率为
决策变量，时段末能量满足
\begin{equation}
E_t^B=(\rho_B)^{\Delta t_t}E_{t-1}^B
+\eta_{ch}P_t^{ch}\Delta t_t
-\frac{P_t^{dis}\Delta t_t}{\eta_{dis}}.
\label{eq:bess-state}
\end{equation}
指数保持结构及初末状态口径与PyPSA公开储能约束一致\cite{pypsa-storage}，电池最优控制
建模可参见Sandia资助研究\cite{sandia-bess}。自放电保持率、充放电效率和可用容量须由
拟选电芯、PCS、温控和运行策略校准。

能量和功率边界为
\begin{equation}
E^{B,r}SOC_{\min}\le E_t^B\le E^{B,r}SOC_{\max},
\label{eq:bess-soc}
\end{equation}
\begin{equation}
0\le P_t^{ch}\le A_t^BP^{ch,\max},\qquad
0\le P_t^{dis}\le A_t^BP^{dis,\max}.
\label{eq:bess-power}
\end{equation}
为禁止同一时段同时充放电，引入二元变量$u_t^{ch},u_t^{dis}$：
\begin{equation}
P_t^{ch}\le u_t^{ch}P^{ch,\max},\quad
P_t^{dis}\le u_t^{dis}P^{dis,\max},\quad
u_t^{ch}+u_t^{dis}\le1.
\label{eq:bess-exclusive}
\end{equation}
电池退化成本按实际吞吐或放电电量计入4.8，具体循环寿命、日历寿命和SOC/温度修正须
来自OEM试验或寿命模型；在获得同一运行边界下的寿命数据前，退化成本系数采用基准值并开展敏感性分析。

\subsection{构网功率与能量备用代理}

若$g_t\in\{0,1\}$表示该时段由BESS承担构网责任，则小时调度需预留向上注入和向下
吸收能力：
\begin{equation}
E^{B,r}SOC_{\min}+g_tE^{GFM,up}
\le E_t^B\le
E^{B,r}SOC_{\max}-g_tE^{GFM,down},
\label{eq:gfm-energy-reserve}
\end{equation}
\begin{equation}
(P_t^{dis}-P_t^{ch})+g_tP^{GFM,up}\le P_t^{dis,\max},
\label{eq:gfm-up-reserve}
\end{equation}
\begin{equation}
(P_t^{ch}-P_t^{dis})+g_tP^{GFM,down}\le P_t^{ch,\max}.
\label{eq:gfm-down-reserve}
\end{equation}
备用功率和能量需依据PCS能力、系统最大扰动及支撑持续时间确定；本轮仅采用基准裕度开展小时级分析。式~\eqref{eq:gfm-energy-reserve}--\eqref{eq:gfm-down-reserve}
只是调度层静态裕度，并非频率最低点、ROCOF、电压恢复、限流、故障穿越或黑启动的
合格证明。相关工作点须进入REGFM、EMT和HIL复核\cite{regfm-b1,unifi-v3}。

\subsection{期末状态与滚动传递}

短窗口方法验证可采用循环边界$E_T^B=E_0^B$，安全储备场景可采用
$E_T^B\ge E^{B,target}$；因果逐时运行不在每个小时强制回到初值，而是将$E_t^B$
原样传递到$t+1$，并通过事件策略、储备下限或经校准的终端价值抑制短视耗尽。期初库存
机会成本用于经济分账，SOC和安全储备仍由物理约束确定。
